\documentclass[11pt,reqno]{amsart}
\usepackage[T1]{fontenc}
\usepackage{lmodern}
\usepackage{microtype}
\usepackage{amsmath,amssymb,mathtools}
\usepackage{booktabs}
\usepackage[colorlinks=true,linkcolor=blue,citecolor=blue,urlcolor=blue]{hyperref}
\newtheorem{theorem}{Theorem}[section]
\newtheorem{lemma}[theorem]{Lemma}
\newtheorem{proposition}[theorem]{Proposition}
\newtheorem{corollary}[theorem]{Corollary}
\newtheorem{introthm}{Theorem}
\renewcommand{\theintrothm}{\Alph{introthm}}
\newtheorem{introcor}[introthm]{Corollary}
\newtheorem{introprop}[introthm]{Proposition}
\theoremstyle{definition}
\newtheorem{definition}[theorem]{Definition}
\newtheorem{remark}[theorem]{Remark}
\newtheorem{assessment}[theorem]{Assessment}
\newtheorem{problem}{Problem}
\newcommand{\Z}{\mathbb Z}
\newcommand{\K}{\mathcal K}
\newcommand{\J}{\mathcal J}
\newcommand{\A}{\mathcal A}
\newcommand{\M}{\mathcal M}
\newcommand{\Sset}{\mathcal S}
\newcommand{\E}{\mathbb E}
\newcommand{\Prob}{\mathbb P}
\newcommand{\ind}{\mathbf 1}
\newcommand{\li}{\operatorname{li}}
\newcommand{\lcm}{\operatorname{lcm}}
\newcommand{\e}{\mathrm e}
\newcommand{\Tabs}{T_{\rm abs}}
\newcommand{\status}[1]{\textup{\textsc{#1}}}
\title[The exponent $3/4$ for $m/n$ in short intervals]{The exponent $3/4$ for the exceptional set of $m/n$, in short intervals and progressions}
\author{Anonymous}
\date{Draft, 8 October 2026; internal checks only, not externally refereed}
\emergencystretch=3em
\begin{document}
\begin{abstract}
For an integer $m\geq4$ call $n$ \emph{$m$-exceptional} if $m/n$ is not a
sum of three unit fractions.  We show that there are absolute constants
$c,C>0$ such that every interval $I$ of length $H\geq2$ contains at most
\[
 C\,H\exp\{-c(\log H)^{3/4}m^{-1/4}\}
\]
$m$-exceptional integers, uniformly in $m\geq4$ and in the position of
$I$.  We also give a version for arithmetic progressions, whose only loss
is the smooth part of the modulus; a bound for exceptional primes in short
intervals that uses no information about primes at the scale of the
interval; and a density statement locating the transition from ``most
primes are $m$-exceptional'' to ``most primes are $m$-representable'' at
$\log n=m^{1/3+o(1)}$, up to a factor $(\log m)^2(m/\varphi(m))^{1/3}$ in
$\log n$.  For comparison, Pomerance and Weingartner proved the
exponent $2/3$ with $\varphi(m)^{-1/3}$ over initial segments.
The method is the one of an earlier internal note on $4/n$ over
$[1,N]$, which we use as a black box.  Two things are new: the transfer
to arbitrary windows is local, through a smooth-part decomposition, and
the numerator $4$ is replaced by $m$, which turns cubic logarithmic class
mass into $(\log X)^3/m$.  \textbf{Status:} every result is proved
\emph{relative to that note}, which is internally proved and has
\emph{not} been externally refereed; the same applies to this note.
Nothing here bears on the existence of exceptions; this is not a proof of
the Erd\H{o}s--Straus conjecture or of Schinzel's conjecture.
\end{abstract}
\maketitle

\section{Introduction}\label{sec:intro}
Let $m\geq2$ be an integer.  A positive integer $n$ is
\emph{$m$-representable} if
\begin{equation}\label{eq:mn}
 \frac mn=\frac1{x_1}+\frac1{x_2}+\frac1{x_3}
\end{equation}
has a solution in positive integers, and \emph{$m$-exceptional} otherwise.
For $m\leq3$ every $n$ is representable; for $m\geq4$ the integer $n=1$
is exceptional.  The Erd\H{o}s--Straus conjecture is the case $m=4$ with
no exception $n\geq2$; Sierpi\'nski conjectured the same for $m=5$, and
Schinzel conjectured that for each $m\geq4$ there are only finitely many
$m$-exceptional $n$.  For a real interval $I$, a modulus $q\geq1$ and
$b\in\Z$ put
\[
 E_m(I)=\#\{n\in I\cap\Z_{\geq1}:n\ m\text{-exceptional}\},\qquad
 E_m(I;q,b)=\#\{n\in I:n\equiv b\ (q),\ n\ m\text{-exceptional}\},
\]
and $E_m(N)=E_m((0,N])$.

\subsection*{Background}
Vaughan \cite{Vaughan} proved $E_4(N)\ll N\exp\{-c(\log N)^{2/3}\}$; this
is the bound quoted as the state of the art by Elsholtz and Tao
\cite[\S1]{ET} and by Pomerance and Weingartner \cite[p.~2]{PW}.
Pomerance and Weingartner \cite[Thm.~1.3]{PW} made it uniform in the
numerator: for $4\leq m\leq(\log N)^2$,
\begin{equation}\label{eq:PW}
 E_m(N)\leq N\exp\{-C_{\rm PW}(\log N)^{2/3}\varphi(m)^{-1/3}\}.
\end{equation}
An internal research note \cite{TQ} (``the note'' below) proves
$E_4(N)\ll N\exp\{-c(\log N)^{3/4}\}$ by an argument that has been
internally checked but not externally refereed.  All results below are
proved relative to \cite{TQ}; Section~\ref{sec:blackbox} lists exactly
which of its statements are used.

\subsection*{Results}
All constants $c,C,\ldots$ below are absolute.  As in \cite{TQ} they
are not effective, because the standard Bombieri--Vinogradov theorem is
used.

\begin{introthm}[short intervals, uniform in $m$]\label{thm:A}
There are absolute constants $c,C>0$ such that for every integer $m\geq4$,
every real $z$ and every $H\geq2$,
\[
 E_m((z,z+H])\leq C\,H\exp\{-c(\log H)^{3/4}m^{-1/4}\}.
\]
In particular $E_m(N)\leq C N\exp\{-c(\log N)^{3/4}m^{-1/4}\}$ for all
$N\geq2$.
\end{introthm}

The bound depends only on the length of the interval.  It contains the
theorem of \cite{TQ}, which is the case $m=4$, $z=0$.  The saving tends
to infinity if and only if $m=o((\log H)^3)$.  For
$m\leq\varepsilon(\log H)^3$ with $\varepsilon=\varepsilon(C)$ small it is
a nontrivial constant factor; for larger $m$ it is no better than the
trivial bound $E_m(I)\leq H+1$.

\begin{introthm}[progressions]\label{thm:B}
There are absolute $c,C,\alpha>0$ with the following property.  Let
$m\geq4$, let $I$ be an interval of length $H$, let $q\geq1$ and $b\in\Z$
with $H/q\geq2$, put
$X_q=\exp\{\alpha(m\log(H/q))^{1/4}\}$, and let $q_1$ be the
$X_q$-smooth part of $q$ (the product of the prime powers $p^a\Vert q$
with $p\leq X_q$).  Then
\[
 E_m(I;q,b)\leq C\,q_1\frac Hq
     \exp\{-c(\log(H/q))^{3/4}m^{-1/4}\}.
\]
Consequently $E_m(I;q,b)\leq C(H/q)\exp\{-\tfrac c2(\log(H/q))^{3/4}m^{-1/4}\}$
whenever $q\leq\exp\{\tfrac c2(\log(H/q))^{3/4}m^{-1/4}\}$, and the
factor $q_1$ is absent when every prime factor of $q$ exceeds $X_q$.
\end{introthm}

\begin{introcor}[exceptional primes in short intervals]\label{cor:C}
With $c,C$ as in Theorem~\ref{thm:A}: for $m\geq4$, $x\geq3$ and $H\geq2$
with $(\log H)^{3/4}m^{-1/4}\geq(2/c)\log\log x$,
\[
 \#\{p\in(x,x+H]:p\text{ prime},\ m\text{-exceptional}\}
 \leq C\frac H{\log x}\exp\{-\tfrac c2(\log H)^{3/4}m^{-1/4}\}.
\]
For fixed $m$ this holds for $H\geq\exp\{C_m(\log\log x)^{4/3}\}$ with
$C_m=(2/c)^{4/3}m^{1/3}$.
\end{introcor}

No information on primes at scale $x$ enters: the factor $1/\log x$ is
paid for by the saving.  Prime-number input is needed only for relative
statements (Remark~\ref{rem:relative}).

\begin{introcor}[the density transition]\label{cor:D}
There is an absolute $C_D$ such that for all $m\geq4$ and all $N\geq16$
with $\log N\geq C_Dm^{1/3}(\log m)^{4/3}$,
\[
 \#\{p\in(N/2,N]:p\text{ prime},\ m\text{-exceptional}\}
 \leq\frac N{(\log N)^2}.
\]
Moreover $E_m(N)=o(N)$ whenever $\log N/m^{1/3}\to\infty$.
\end{introcor}

In the other direction, the proof of \cite[Thm.~3.1]{PW} shows that most
primes in $(N/2,N]$ are $m$-exceptional when
$m^{1/4}\ll\log N\leq(\varphi(m)/(C\log^2m))^{1/3}$
(Section~\ref{sec:transition}).  Pomerance and Weingartner's rigorous
bounds placed the transition between $\log N=m^{1/3}$ and
$\log N=m^{1/2}$ \cite[p.~2]{PW}; their Poisson heuristic \cite[p.~3]{PW}
predicts $m^{1/3+o(1)}$.  Corollary~\ref{cor:D} confirms the density
part of that prediction, up to the factor
$(\log m)^2(m/\varphi(m))^{1/3}$ in $\log N$.  We do not know that the
proportion of exceptional primes is monotone in between, and nothing is
said about the \emph{largest} $m$-exceptional integer.

Theorem~\ref{thm:A} and \eqref{eq:PW} have unrelated ineffective
constants, so their comparison is asymptotic.  The ratio of the
exponents is
\begin{equation}\label{eq:ratio}
 \frac{L^{3/4}m^{-1/4}}{L^{2/3}\varphi(m)^{-1/3}}
 =L^{1/12}\frac{\varphi(m)^{1/3}}{m^{1/4}}
 \gg\frac{(Lm)^{1/12}}{(\log\log3m)^{1/3}},\qquad L=\log N,
\end{equation}
by $\varphi(m)\gg m/\log\log3m$.  Hence there is an absolute
(ineffective) $\Lambda_0$ such that Theorem~\ref{thm:A} is stronger than
\eqref{eq:PW} whenever $(Lm)^{1/12}\geq\Lambda_0(\log\log3m)^{1/3}$; in
particular for all $4\leq m\leq L^2$ once $N\geq N_0$.

Finally (Section~\ref{sec:lowerend}) a bound with saving in $\log x$
rather than $\log H$, of the shape
$E_m((x,x+H])\ll H\e^{-c(\log x)^{3/4}}$, cannot be proved for windows
with $H<\e^{c(\log x)^{3/4}}/C$ without proving that there are no
exceptions at all beyond some point (Proposition~\ref{prop:E}).

\subsection*{Status and novelty}
Every statement labelled \status{proved} here is proved \emph{relative to}
\cite{TQ}: the steps that change are written out in full, and the
remaining steps are cited from \cite{TQ} by name.  The note \cite{TQ}
is internally proved; it has \emph{not} been refereed externally, and
neither has this note.  The underlying working documents
\cite{SHORT,MN} received three internal hostile reviews, and their
repairs are incorporated.

We found no published short-interval or progression version of an
exceptional-set bound for \eqref{eq:mn}, for any $m$
(Section~\ref{sec:literature}); the searches were not exhaustive.  We
make no novelty claim for the \emph{method}, which is that of \cite{TQ},
and the forced classes are those of Vaughan, Elsholtz--Tao and
Pomerance--Weingartner.  Every forced class is an identity valid for all
$n$ in the class, and counting congruence classes is translation
invariant.  An expert would therefore regard the transfer to short
intervals as routine for the part of the argument that concerns integers
free of small prime factors.  What we add is the following: a local
replacement for the global semigroup step of \cite{TQ}
(Lemma~\ref{lem:smooth}); the bookkeeping that turns $t^3$ into $t^3/m$
with absolute constants, with no loss of $\varphi(m)$ or $\log m$
(Sections~\ref{sec:atoms}--\ref{sec:assembly}); and the consequences
above.  None of this improves the exponent $3/4$.

\subsection*{Organisation}
Section~\ref{sec:blackbox} lists what is taken from \cite{TQ} and
isolates the single object the rest of the argument needs, a
\emph{majorant package} (Definition~\ref{def:package}).
Sections~\ref{sec:atoms}--\ref{sec:assembly} construct the package for
general $m$, writing out every step that differs from \cite{TQ}.
Section~\ref{sec:transfer} contains the local transfer to arbitrary
windows and progressions and proves Theorems~\ref{thm:A},~\ref{thm:B}
and Corollary~\ref{cor:C}.  Section~\ref{sec:transition} proves
Corollary~\ref{cor:D}, Section~\ref{sec:lowerend} discusses the lower
end of the window range, Section~\ref{sec:literature} the literature,
and Section~\ref{sec:open} lists open problems.

\section{The black box and the majorant package}\label{sec:blackbox}
\subsection{Parameters}
We use the parameters of \cite[(2)]{TQ} with one change.  Fix
$\kappa=1/480$ and the integer $D$ of \cite[Lemma~3.3]{TQ}.  For $X\geq3$
put
\[
 t=\log X,\qquad K=\lfloor X^\kappa\rfloor,\qquad H_K=K^{10},
\]
and use the blocks $I_j=(x_j,\min(2x_j,X)]$, $x_j=2^jX^{1/2}$, with
$z_j=x_j^{1/6}$; for large $X$ one has $z_j\geq X^{1/12}>H_K^2$
\cite[(3)]{TQ}.  (The note calls $H_K$ simply $H$; here $H$ is the length of
a window.)  The change is the multiplier set:
\[
 \K_m(K)=\{1\leq k\leq K:(k,m)=1\},\qquad L_K=\lcm\K_m(K),
\]
in place of $\K(K)=\{k\leq K:k\equiv1\ (4)\}$.  For
$\J\subseteq\K_m(K)$ put $h(\J)=\sum_{k\in\J}\varphi(k)/k^2$ and
$L_{\J}=\lcm\J$.  Throughout, $P_y=\prod_{p\leq y}p$ and
$S_y(n)=\ind_{(n,P_y)=1}$.

\subsection{What is used from the note}
External inputs are those of \cite[\S3]{TQ}: Brun--Titchmarsh in the
Montgomery--Vaughan form \cite{MV}, the Bombieri--Vinogradov theorem
\cite[Ch.~28]{Davenport}, and Shiu's theorem \cite{Shiu} (inside
\cite[Lemma~3.3]{TQ}).  The statements of \cite{TQ} enter as follows.

\smallskip
\noindent\textbf{Used verbatim.}
\begin{itemize}
\item \cite[Lemma~3.1]{TQ} (elementary lattice sums); it is stated for
every integer $1\leq k\leq K$ and every $(c,k)=1$.
\item \cite[Lemma~3.3]{TQ} (low-$\omega$ mass), stated for every
$k\leq K$, $(c,k)=1$.
\item \cite[Lemma~4.1]{TQ} (second incidence moment).  It is stated for
$\J\subseteq\K(K)$, but its proof uses only $\J\subseteq[1,K]$, so it
holds for $\J\subseteq\K_m(K)$.
\item The conditional-independence proof of \cite[Theorem~6.3]{TQ}
(factorial moments) and \cite[Lemma~8.1]{TQ} (the even Bonferroni
polynomial $Q_r$).
\item For $m=4$ only: \cite[Theorem~8.2]{TQ} (the assembly) as stated,
together with the identity \cite[Lemma~2.1]{TQ}.
\end{itemize}

\noindent\textbf{Re-run with $4$ replaced by $m$} (each change is written
out below): \cite[Lemma~2.1]{TQ} (the identity; Lemma~\ref{lem:identity}),
\cite[Lemma~2.2]{TQ} (deduplication and CRT; Lemma~\ref{lem:CRT}),
\cite[Lemma~3.2]{TQ} (harmonic sum; replaced by Lemma~\ref{lem:hm}),
\cite[Theorem~4.2, Corollary~4.3]{TQ} (prime slice and fibre masses;
Proposition~\ref{prop:harvest}, Corollary~\ref{cor:fibre}),
\cite[Lemma~7.1]{TQ} (void; Lemma~\ref{lem:void}) and
\cite[Theorem~8.2]{TQ} (assembly; Theorem~\ref{thm:assembly}).

\noindent\textbf{Replaced.}  The counting step \cite[(50)]{TQ} over
$[1,N]$ is replaced by its local form (Lemma~\ref{lem:localmean}), and
the global transfer to all denominators \cite[\S9]{TQ} (semigroup of
exceptional primes and Rankin's trick) is replaced by a local smooth-part
decomposition (Lemma~\ref{lem:smooth}).

\noindent\textbf{Not used.}  The simplified route \cite[\S5]{TQ}, the
ordered-atom replay after \cite[Theorem~6.3]{TQ} with
\cite[Lemmas~6.1, 6.2]{TQ}, and \cite[\S\S10--11]{TQ} (heuristic ceiling,
checkpoints).  We re-derive only the retained original route of
\cite{TQ}, with the cutoff $\omega(uv)\leq D\log\log X$ and congestion
pruning.  (The Cauchy--Schwarz route of \cite[\S5]{TQ} also appears to
transfer, via $\varphi(muv)\geq\varphi(m)\varphi(u)\varphi(v)$, but we
have not written it out.)

\smallskip
The substitutions are summarised in Table~\ref{tab:subst}.  Where the
numerator $4$ occurs in \cite{TQ} we checked it occurrence by occurrence.
It appears only in the identity, in the multiplier condition
$k\equiv1\ (4)$ and \cite[Lemma~3.2]{TQ}, in the prime modulus $4uv$, in
the inequality $\varphi(4uv)\geq2\varphi(u)\varphi(v)$, and in the
non-load-bearing \cite[\S\S10--11]{TQ}.  The inventory bound
$|\A_X|\leq KX^{4/3}$ does not involve the numerator.  No step of the
density argument (lattice sums, Shiu, congestion, moments, void,
Bonferroni) looks at the quadratic character of the forced classes.  So
the dichotomy between $m\equiv0\pmod 4$ and $m\not\equiv0\pmod4$ that
governs \emph{pointwise} questions for $m/n$ \cite[\S15]{SUBEXP} plays no
role here.

\begin{table}[ht]
\centering
\small
\begin{tabular}{@{}lll@{}}
\toprule
object & note \cite{TQ} ($m=4$) & here (general $m$)\\
\midrule
identity & $k\ell+1=4uvw$ & $k\ell+1=m\,uvw$\\
multipliers & $k\leq K$, $k\equiv1\ (4)$ & $k\leq K$, $(k,m)=1$\\
prime progression & $\ell\equiv-k^{-1}\ (4uv)$ & $\ell\equiv-k^{-1}\ (muv)$\\
totient bound & $\varphi(4uv)\geq2\varphi(u)\varphi(v)$ &
   $\varphi(muv)\geq\varphi(m)\varphi(u)\varphi(v)$\\
harmonic sum & $h=\tfrac2{\pi^2}\log K+O(1)$ &
   $h_m\asymp(\varphi(m)/m)\log K$\\
fibre mass & $\asymp t^3$ & $\asymp s:=t^3/m$\\
selector, depth & $y=Bt^3$, $r\approx D_Bt^3$ & $y=\max(2,Bs)$, $r\approx D_Bs$\\
ledger & $\e^{O(t^4)}$ & $\e^{O(ts)}=\e^{O(t^4/m)}$\\
transfer & $[1,N]$, then Rankin & any window, smooth parts\\
\bottomrule
\end{tabular}
\caption{The substitutions $4\to m$ and $[1,N]\to$ window.}
\label{tab:subst}
\end{table}

\subsection{The majorant package}
The transfer of Section~\ref{sec:transfer} uses only the following
properties.

\begin{definition}\label{def:package}
Let $m\geq4$, $X\geq3$, $t=\log X$, $s\geq1$ and $2\leq y<X^{1/2}$.  An
\emph{$(m;X,s,y)$-package with constants $(a,C_L)$} is a function
$\nu:\Z\to\mathbb R$ together with a positive integer $\M$ all of whose
prime factors are at most $X$, such that:
\begin{enumerate}
\item[(P1)] $\nu(n)\geq0$ for all $n$, and $\nu(n)\geq1$ for every
$m$-exceptional $n\geq1$ with $(n,P_y)=1$;
\item[(P2)] $\nu=\sum_i\epsilon_i\ind_{a_i\ (q_i)}$ is a finite signed
combination of residue classes with every $q_i\mid\M$ and
$\Tabs:=\sum_i|\epsilon_i|\leq\e^{C_Lts}$;
\item[(P3)] $\E_{\M}\nu\leq2\e^{-as}$, where $\E_{\M}$ is the uniform
average over $\Z/\M\Z$ (well defined by (P2)).
\end{enumerate}
\end{definition}

\begin{proposition}[packages exist]\label{prop:package}
There are absolute constants $a,B,C_L,s_0>0$ and $X_a$ such that for
every $X\geq X_a$ and every integer $4\leq m\leq t^3$ with
$s:=t^3/m\geq s_0$ there is an $(m;X,s,\max(2,Bs))$-package with
constants $(a,C_L)$.
\end{proposition}

For $m=4$ the note gives such a package directly.  Indeed
\cite[Theorem~8.2]{TQ}, with its own family and fixed constants, yields a
$(4;X,t^3/4,B_4t^3)$-package with constants $(4c_a,4C_L)$ in the notation
of \cite{TQ}: property (P1) holds because every
exceptional $n$ avoids every atom \cite[Lemma~2.2]{TQ}, so $H_X(n)=0$ and
$\nu_X(n)=S_y(n)Q_r(0)=1$ when $(n,P_y)=1$; (P2) is \cite[(44)]{TQ} with
$\M$ from \cite[(7)]{TQ}; and (P3) is \cite[(43)]{TQ}.  Note that
\cite[Theorem~8.2]{TQ} states the lower bound $\nu_X\geq1$ only for
exceptional primes $n>\max(K,y)$.  Its proof, however, uses only the
identity, which holds for every positive integer in a forced class, so
it gives (P1) as stated.  The transfer of Section~\ref{sec:transfer}
needs only $y\leq B's+2$ for some absolute $B'$, which holds here with
$B'=4B_4$.  Thus the case $m=4$ of Theorems~\ref{thm:A} and~\ref{thm:B}
follows from \cite{TQ} and Section~\ref{sec:transfer} alone.  For general $m$,
Proposition~\ref{prop:package} is Theorem~\ref{thm:assembly} below.

\begin{thebibliography}{99}
\bibitem[BHP]{BHP}
R. C. Baker, G. Harman and J. Pintz, The difference between consecutive
primes, II, \emph{Proc. London Math. Soc. (3)} \textbf{83} (2001),
532--562.
\bibitem[D]{Davenport}
H. Davenport, \emph{Multiplicative Number Theory}, third edition,
revised by H. L. Montgomery, Graduate Texts in Mathematics 74,
Springer, 2000.
\bibitem[ET]{ET}
C. Elsholtz and T. Tao, Counting the number of solutions to the
Erd\H{o}s--Straus equation on unit fractions, \emph{J. Aust. Math.
Soc.} \textbf{94} (2013), 50--105; arXiv:1107.1010.
\bibitem[MN]{MN}
\emph{EXCEPTIONAL\_MN: the $3/4$ exceptional-set bound for $m/n$},
internal working document \texttt{EXCEPTIONAL\_MN.md}, this repository,
2026; hostile reviews \texttt{reviews/exceptional-mn-review-A.md},
\texttt{reviews/exceptional-mn-review-B.md}.
\bibitem[MV]{MV}
H. L. Montgomery and R. C. Vaughan, The large sieve,
\emph{Mathematika} \textbf{20} (1973), 119--134.
\bibitem[PW]{PW}
C. Pomerance and A. Weingartner, Exceptions to the
Erd\H{o}s--Straus--Schinzel conjecture, arXiv:2511.16817v2, 2026.
\bibitem[S]{Shiu}
P. Shiu, A Brun--Titchmarsh theorem for multiplicative functions,
\emph{J. Reine Angew. Math.} \textbf{313} (1980), 161--170.
\bibitem[SH]{SHORT}
\emph{EXCEPTIONAL\_SHORT: exceptional sets in short intervals and
progressions}, internal working document \texttt{EXCEPTIONAL\_SHORT.md},
this repository, 2026; hostile review
\texttt{reviews/exceptional-short-review.md}.
\bibitem[SX]{SUBEXP}
\emph{Large Erd\H{o}s--Straus witnesses II}, research note,
\texttt{paper/es-subexp-note.tex}, this repository, 2026.
\bibitem[TQ]{TQ}
\emph{A three-quarter exceptional-set bound for the Erd\H{o}s--Straus
equation: an internally proved research note},
\texttt{paper/es-threequarter-note.tex}, this repository, 2026;
internally proved, not externally refereed.  Numbering of statements
and displays as in the compiled version.
\bibitem[V]{Vaughan}
R. C. Vaughan, On a problem of Erd\H{o}s, Straus and Schinzel,
\emph{Mathematika} \textbf{17} (1970), 193--198.
\end{thebibliography}
\end{document}
